\documentclass[aspectratio=169]{beamer}
\usetheme{Madrid}
\usecolortheme{default}
\usepackage[utf8]{inputenc}
\usepackage{booktabs}
\usepackage{amsmath}

\title{Ante-hoc Explainable AI for\\Foreshock vs Aftershock Classification}
\author{XAI-CH Project}
\date{\today}

\begin{document}

\begin{frame}
  \titlepage
\end{frame}

\begin{frame}{Task Overview: Exploring an Ante-hoc XAI Method}
\begin{itemize}
  \item Goal: classify seismic events as \textbf{foreshock} (label 0) or \textbf{aftershock} (label 1), while keeping model reasoning interpretable.
  \item Ante-hoc explainability means the model is designed to be interpretable \emph{during} training/inference, not explained only post-hoc.
  \item Repository notebooks evaluate two approaches:
  \begin{itemize}
      \item \textbf{Explainable Boosting Machine (EBM)} on engineered 3-segment features.
      \item \textbf{Token-based Attention Transformer} on tokenized waveform segments.
  \end{itemize}
  \item Dataset size reported in notebooks: \textbf{2,341} foreshock/aftershock events.
\end{itemize}
\end{frame}

\begin{frame}{Method 1: EBM (Feature-based Ante-hoc Model)}
\textbf{How the model works}
\begin{itemize}
  \item EBM is a generalized additive model (GAM) with boosted trees for each feature (and optionally selected interactions).
  \item Prediction is additive:
  \[
    \hat{y} = b + \sum_j f_j(x_j)\ (+\ \sum_{j,k} f_{jk}(x_j,x_k)\,\text{if interactions are enabled})
  \]
  \item Notebook uses 3-segment seismic features (noise, P-phase, coda) and channel-wise spectral/statistical descriptors.
\end{itemize}

\textbf{Why it is explainable}
\begin{itemize}
  \item Global explainability: feature importance and per-feature shape functions.
  \item Local explainability: event-level additive contribution of each feature.
  \item Additive structure supports direct traceability from feature values to decision score.
\end{itemize}
\end{frame}

\begin{frame}{EBM Results from Notebook}
\begin{itemize}
  \item Test-set metrics reported:
  \begin{itemize}
    \item Accuracy: \textbf{0.6077}
    \item F1-score: \textbf{0.6290}
    \item ROC-AUC: \textbf{0.6364}
  \end{itemize}
  \item Class-wise report highlights:
  \begin{itemize}
    \item Foreshock: precision 0.62, recall 0.55, F1 0.58.
    \item Aftershock: precision 0.60, recall 0.66, F1 0.63.
  \end{itemize}
  \item Confusion matrix indicates moderate but meaningful separability between classes.
\end{itemize}
\end{frame}

\begin{frame}{EBM Explanation Examples from Notebook}
\textbf{Global feature importance examples}
\begin{itemize}
  \item Top features include:
  \begin{itemize}
    \item \texttt{N\_coda\_band\_20\_45\_pow} (importance 0.1023)
    \item \texttt{N\_coda\_band\_5\_10\_pow} (importance 0.0987)
    \item \texttt{N\_noise\_band\_5\_10\_pow} (importance 0.0701)
    \item \texttt{Z\_noise\_band\_20\_45\_pow} (importance 0.0666)
  \end{itemize}
  \item Segment/channel aggregation in notebook ranks \textbf{coda-N} and \textbf{coda-Z} among strongest contributors.
\end{itemize}

\textbf{Interpretation result}
\begin{itemize}
  \item The model emphasizes coda/noise spectral-energy patterns, providing a physically interpretable pathway for class separation.
\end{itemize}
\end{frame}

\begin{frame}{Method 2: Token-based Attention Transformer}
\textbf{How the model works}
\begin{itemize}
  \item Waveforms are split into fixed-length tokens; each token is projected to an embedding.
  \item A Transformer encoder with multi-head self-attention processes token sequence plus a CLS token.
  \item Final class probability is produced from the CLS representation.
\end{itemize}

\textbf{Why it can be explainable}
\begin{itemize}
  \item Last-layer CLS$\rightarrow$token attention provides token-level relevance scores.
  \item Relevance can be mapped back to event timeline to inspect what temporal regions influence the decision.
  \item Notebook additionally evaluates faithfulness (removal/insertion curves), testing if high-attention tokens are truly influential.
\end{itemize}
\end{frame}

\begin{frame}{Transformer Results from Notebook}
\begin{itemize}
  \item Test-set metrics reported:
  \begin{itemize}
    \item Accuracy: \textbf{0.4886}
    \item F1-score: \textbf{0.4479}
    \item ROC-AUC: \textbf{0.4926}
  \end{itemize}
  \item Class-wise report highlights:
  \begin{itemize}
    \item Foreshock: precision 0.49, recall 0.56, F1 0.52.
    \item Aftershock: precision 0.49, recall 0.41, F1 0.45.
  \end{itemize}
  \item Performance is near chance-level in this run, lower than EBM in all three metrics.
\end{itemize}
\end{frame}

\begin{frame}{Transformer Explanation Examples from Notebook}
\textbf{Attention-based examples}
\begin{itemize}
  \item Notebook visualizes last-layer CLS attention over tokens aligned with waveform timeline.
  \item This provides event-level explanation: which temporal tokens were weighted most for each prediction.
\end{itemize}

\textbf{Faithfulness test results}
\begin{itemize}
  \item Base mean $P(\text{class}=1)$: \textbf{0.4898}
  \item Removal gap at 50\% tokens kept (attention-random): \textbf{0.2305}
  \item Insertion gap at 50\% tokens inserted (attention-random): \textbf{-0.0612}
\end{itemize}

\textbf{Interpretation result}
\begin{itemize}
  \item Attention carries some signal for removal-based perturbation, but insertion result suggests incomplete faithfulness/calibration.
\end{itemize}
\end{frame}

\begin{frame}{Overall Comparison and Takeaways}
\begin{table}
\centering
\begin{tabular}{lccc}
\toprule
Method & Accuracy & F1 & ROC-AUC \\
\midrule
EBM (feature-based) & 0.6077 & 0.6290 & 0.6364 \\
Token Attention Transformer & 0.4886 & 0.4479 & 0.4926 \\
\bottomrule
\end{tabular}
\end{table}

\vspace{0.4cm}
\begin{itemize}
  \item In these notebook runs, \textbf{EBM is both more accurate and more reliably interpretable}.
  \item Transformer gives richer temporal explanations, but predictive quality and faithfulness need improvement.
  \item Practical next step: preserve token-level explainability while improving architecture/training and validating with stronger faithfulness checks.
\end{itemize}
\end{frame}

\begin{frame}{End}
\centering
\Large Questions?
\end{frame}

\end{document}
